% !TEX root = ../Thesis.tex
% ============================================================
\chapter{Methodology} \label{chapter:methodology}
% ============================================================

Every choice described below is a decision with a measured justification and an
off switch. That constraint is deliberate. A pipeline whose steps cannot be
disabled cannot be ablated, and a step that cannot be ablated cannot honestly be
reported as a contribution, because its effect stays indistinguishable from the
effect of the pipeline as a whole. Decisions that remain open are marked as open
rather than quietly settled.

\section{Graph Preparation} \label{sec:methodology-preparation}

An SPHN instance graph is not directly amenable to text-based retrieval, for
reasons established in Section~\ref{subsec:sphn-schema}: descriptions sit on
shared code nodes rather than on the events a question is about, and provenance
nodes of enormous degree make connectivity vacuous. Four transformations address
this. Each is applied to a derived text layer; none modifies the database, which
is read-only by design.

\paragraph{Patient scoping.} Retrieval is performed within one patient's
subgraph. This is not merely an optimisation: a question about a patient has no
answer outside that patient, and scoping removes $1{,}196/1{,}197$ of the graph
before anything else runs. Cohort-level questions are out of scope for the metric
for a separate reason discussed in Section~\ref{sec:methodology-metric}.

\paragraph{Folding shared terminology into text.} Nodes whose only role is to
carry a value or a code (\texttt{Code}, \texttt{Unit}, \texttt{Quantity},
\texttt{BodySite} and \texttt{ReferenceRange}) are removed as nodes, their text
composed into the description of the event that references them. This is what
turns a four-node laboratory construct plus its satellites into one node carrying
a sentence. Without it, the node a question is about carries a timestamp and
nothing else.

\paragraph{Pruning provenance.} \texttt{SourceSystem},
\texttt{HealthcarePrimaryInformationSystem}, \texttt{DataProvider},
\texttt{DataRelease}, \texttt{TimePattern}, \texttt{CareHandling} and
\texttt{Location} are removed. They are correct, schema-mandated, and place
almost every pair of nodes within three hops of each other
(Table~\ref{tab:measured-graph}). The justification is not size but
interpretability of connectivity: a path through a provenance hub carries no
clinical meaning.

\paragraph{Two text fields per node.} Each node is given an \emph{embedding text}
and a \emph{display text}. The embedding text carries the semantic core, meaning
what kind of measurement and which analyte, and deliberately excludes the measured
value and the timestamp, so that repeated measurements of one analyte share a
single embedding. The display text carries everything, because an LLM answering
from the subgraph needs the value. The split is what makes embedding a
15{,}810-node patient cost 614 encoder calls rather than 15{,}810, and it is also
the reason a text-scoring baseline must be given the embedding text and not the
display text if the comparison is to be fair.

\paragraph{Terminology resolution.} Codes are resolved against published
catalogues to supply descriptions the graph lacks, across five tiers: exact
match, a different version of the same classification, the parent concept, the
graph's own label, and nothing. The count of nodes answered by each tier is
reported alongside any result depending on it. The design is a cross-walk rather than a
translation because a code is language-independent, which makes the operation
deterministic and auditable, and because it can fill codes that have no label at
all. Section~\ref{sec:results-terminology} reports what it recovers.

Each of the five is independently switchable, and the un-prepared graph is the
baseline every gain is measured against.

\section{Two-Stage Retrieval} \label{sec:methodology-candidates}

G-Retriever assumes one modest graph per question. A prepared patient is
$15{,}810$ nodes (Table~\ref{tab:patient-zero}), so a narrowing stage is
introduced before selection: a generator proposes a region, and the selection
strategy under test operates inside it. Two generators are compared, top-$n$ by
similarity and expansion from similarity-ranked seeds under a hard cap, together
with the control of no narrowing at all.

Narrowing can discard an answer before any retriever sees it, and that failure is
indistinguishable from a bad selection if only the final subgraph is scored. Every
two-stage result is therefore reported with a \emph{ceiling}: the recall
attainable given what stage one kept. Recall far below its ceiling indicates a
selection problem; a low ceiling indicates a narrowing problem.

\section{Strategies Compared} \label{sec:methodology-strategies}

Five strategies are evaluated, one per family of
Section~\ref{sec:subgraph-selection-approaches}: PCST-based selection; top-$k$
triples in the manner of KAPING~\cite{baek2023kaping}; top-$k$ nodes with their
neighbours; breadth-first expansion from seeds; and shortest paths between
similarity-ranked nodes. All five consume the same graph, the same encoder and
the same question embedding, and differ only in how they choose. Any baseline
that scores text rather than embeddings is given the embedding text, since
scoring it on display text while scoring PCST on embedding text would understate
it, and understating a baseline is as much an error as overstating the method
under test.

\section{The Metric} \label{sec:methodology-metric}

The primary measurement is \emph{answer-node recall at a given subgraph size}:
the fraction of the entities constituting a correct answer that appear in the
retrieved subgraph, reported against the measured size of that subgraph.

For a question with gold answer nodes $A$ and a retrieved node set $V_S$,
%
\begin{equation}
\mathrm{recall} = \frac{\lvert A \cap V_S \rvert}{\lvert A \rvert} .
\label{eq:recall}
\end{equation}
%
Recall is undefined when $A$ is empty and is recorded as such rather than as
zero. This matters more than it appears: a question such as ``how many patients
received drug X'' has no answer node anywhere in the graph, and scoring it zero
would report a well-behaved retriever as failing. Such questions are recorded
with a reason and excluded from the mean, and the number of questions actually
scored is reported beside every average.

Alongside recall, each run records the retrieved node count, edge count, prompt
character count, the fraction of the graph retrieved, and whether the result is
connected.

\section{Comparing Dials That Are Not the Same Quantity} \label{sec:methodology-comparability}

PCST's size parameter is a cost per edge; the baselines' is a count $k$ or a hop
number $h$. These are not commensurable, and plotting recall against ``the
parameter'' would compare a cost with a count. Every comparison in this thesis is
therefore made against a \emph{measured} size, either node count or prompt
characters, with each strategy swept across its own parameter range and the
resulting points placed on a shared axis.

Two consequences must be stated rather than assumed. First, a strategy that
returns more nodes will usually recall more, so a recall figure is meaningless
without the size beside it; Section~\ref{sec:results-candidates} contains a case
where ignoring this would have manufactured an improvement. Second, the ranges
must overlap in achieved size, or the curves never meet and no comparison at
equal size is possible.

\section{Question Sets} \label{sec:methodology-questions}

Three sets are used, with decreasing availability and increasing authority.

\paragraph{Planted questions.} Generated against a synthetic graph with known
answer nodes. They test the harness and structural mechanisms and make no claim
about clinical retrieval. They are the only set available without access to real
usage, and every result in Chapter~\ref{chapter:results} rests on them.

\paragraph{Real user questions.} Questions clinicians and data managers have
actually put to SnapQuery. Whether these are retained, in what form, and who can
export them is an open question with the STCS Data Center at the time of writing.
Their availability determines whether RQ1 can be answered on real usage or only
on a curated set.

\paragraph{A validated gold set.} Approximately fifty questions whose answer
entities have been confirmed by hand, so that correctness is established
independently of SnapQuery's own output. Using a production system's answers as
ground truth for a comparison against that same system would be circular; the
gold set is what breaks the circle, which is why it is treated as a core artefact
rather than a validation afterthought.

\subsection{Recording a Gold Question} \label{sec:methodology-goldset}

A gold question carries its text, the language it is written in, whether its
answer is a set of entities or an aggregate, the patient it is scoped to, its
answer entities, and who confirmed them on what date. Language is recorded
rather than inferred because it decides both the encoder and the terminology
catalogue, and Section~\ref{sec:results-terminology} shows what applying the
wrong one costs.

Answers are stored as \emph{stable keys} derived from the SPHN uid, never as
node indices. An index is an artefact of the order in which the loader happened
to walk the database, and it moves whenever the loader, a query limit or the
data changes. An answer recorded as index $4{,}812$ would go on scoring after it
had begun denoting a different node, and no error would be raised. Resolution
against a graph therefore reports keys it cannot find instead of quietly
skipping them, since a missing key means the loader no longer emits a node the
gold set was built against, and silently scoring around it would understate
every method equally and invisibly.

Four conditions are refused outright, each corresponding to a way a question set
measures something other than retrieval. An entity question with no answers
scores zero against every method. An aggregate question that carries answers
gets scored as though recall meant something for it. A question that contains
the code it is looking for is answerable by string overlap alone; this defect
was found in the planted-question generator, where it made PCST appear perfect.
A question with no named validator has another system's output as its ground
truth, which is the circularity the set exists to avoid, and is reported
separately rather than rejected.

\section{Evaluating Across Patients} \label{sec:methodology-multipatient}

Retrieval is scoped to one patient, so a cohort result is an aggregate over
per-patient runs rather than a single measurement. Three properties of that loop
matter more than the retrieval it performs.

Failures are outcomes. Loading a patient can fail for reasons unconnected to
retrieval, and a run over $1{,}197$ patients that aborts on the fortieth has
produced nothing. Each failure is caught, recorded with its reason, and reported
as a count beside the results.

Runs are resumable. Loading and encoding dominate the cost, so rows are appended
as they are produced and finished patients are skipped on restart.

Patients are averaged, not questions. Pooling every question into one mean lets
a few large patients dominate and conceals the variance, which is itself a
result. Each patient contributes one value to the cohort mean, and the standard
deviation across patients is reported beside it.

\section{SnapQuery as a Comparator} \label{sec:methodology-snapquery}

SnapQuery returns rows, not a subgraph (Section~\ref{subsec:snapquery-background}).
To place it on the same axis, the entities its result rows touch are treated as
its induced subgraph. Two limitations follow and are stated here rather than
discovered later. It exposes no size parameter, so it contributes a single point
rather than a curve, and comparisons must be read as ``at the size SnapQuery
happens to produce''. And the mapping from rows to entities is an interpretation
imposed by this thesis, not something SnapQuery asserts; a row that aggregates
over entities may touch many or none.

\section{Data Governance} \label{sec:methodology-governance}

The instance data is governed clinical data and remains on University Hospital
Basel and BioMedIT infrastructure throughout. All database access is read-only.
Only schema names, property names, counts and aggregate statistics appear in this
thesis; no property value, identifier, laboratory result or date does. Patients
are addressed by position in a query result rather than by any identifier, so no
pseudo-identifier leaves the server. Development proceeds against the public
schema, against synthetic graphs reproducing its measured structure, and against
the real data only for the measurements reported as such.
